(1) 横向力平衡给 $\mu y_{tt}=Ty_{xx}-2\mu\gamma y_t$，即 $y_{tt}+2\gamma y_t=c^2y_{xx}$，$c=\sqrt{T/\mu}$。固定端 $y(0,t)=y(L,t)=0$ 使空间本征模为 $\sin(n\pi x/L)$、波数 $k_n=n\pi/L$，$n=1,2,\ldots$。
(2) 设 $y(x,t)=q_n(t)\sin k_nx$，可得 $\ddot q_n+2\gamma\dot q_n+\omega_{0n}^2q_n=0$，$\omega_{0n}=ck_n$。当 $\gamma<\omega_{0n}$ 时 $q_n=e^{-\gamma t}[A_n\cos(\omega_dt)+B_n\sin(\omega_dt)]$，$\omega_d=\sqrt{\omega_{0n}^2-\gamma^2}$，振幅包络确为 $e^{-\gamma t}$。
(3) 能量 $E=\frac12\int_0^L[\mu y_t^2+T y_x^2]dx$。微分并分部积分，固定端不做功，得 $dE/dt=-2\mu\gamma\int_0^L y_t^2dx\le0$。弱阻尼下周期平均动能约占总能量一半，故能量包络为 $e^{-2\gamma t}$，每弧度耗散比给 $Q_n\simeq\omega_{0n}/(2\gamma)$。
(4) 模态满足 $\gamma<ck_n$ 才振荡，等号临界阻尼，$\gamma>ck_n$ 时为两个实衰减根的非振荡解。$\gamma\to0$ 恢复标准弦的正弦时间振动且能量守恒；强阻尼不能把 $\omega_d$ 当成实数振荡频率，也不应继续套用弱阻尼品质因数近似。
